% =====================================================================
%  4.3 Derivasjonsregler for polynom — teori og oppgaver
% =====================================================================

\del{1}{Derivere polynomer}

\begin{ordliste}
derivasjonsregel & En regel som gir den deriverte uten å gå via grenseverdi. & \mittsprak\\
polynom & En sum av ledd på formen $ax^n$, for eksempel $2x^3 - 5x + 4$. & \mittsprak\\
fart & Hvor fort posisjonen endrer seg. Farten er den deriverte av strekningen. & \mittsprak\\
momentan vekstfart & Vekstfarten i ett bestemt punkt. Det er $f'(a)$. & \mittsprak\\
\end{ordliste}

\begin{mustkunne}
\begin{minipage}[t]{0.50\linewidth}
\textbf{Reglene fra 4.2 samlet}\par
\renewcommand{\arraystretch}{1.4}
\begin{tabular}[t]{@{}l l@{}}
$(k)' = 0$ & en konstant vokser ikke\\
$(kx)' = k$ & stigningstall $k$\\
$(x^n)' = n\,x^{n-1}$ & mønsteret fra 4.2\\
$(k\cdot f)' = k\cdot f'$ & konstant utenfor\\
$(f + g)' = f' + g'$ & vekstfartene legges sammen\\
\end{tabular}
\end{minipage}\hfill
\begin{minipage}[t]{0.47\linewidth}
\textbf{Grensen i 1T.} Vi har ingen regel for $f\cdot g$ eller for brøker. Derfor gjør vi alltid dette først:
\begin{itemize}[leftmargin=1.2em,itemsep=1pt,topsep=2pt]
  \item gang ut produkter, for eksempel $(x+1)^2$
  \item forkort brøker, for eksempel $\frac{x^3 - 2x}{x} = x^2 - 2$ ($x \neq 0$)
\end{itemize}
\textbf{Tolkning:} $s'(t)$ er farten, $K'(x)$ er omtrent kostnaden for én enhet til. Skriv alltid enhet.
\end{minipage}
\end{mustkunne}

\begin{eksempel}[1 — derivér og tolk]
La $f(x) = 2x^3 - 5x^2 + 4$. Finn $f'(x)$ og $f'(2)$.

\begin{stegtabell}
$f'(x) = 2\cdot 3x^2 - 5\cdot 2x + 0$ & Hvert ledd for seg: $(x^n)'$, så konstant faktor. Konstanten $4$ gir $0$.\\
$f'(x) = 6x^2 - 10x$ & Regn ut.\\
$f'(2) = 6\cdot 2^2 - 10\cdot 2 = 24 - 20 = 4$ & Sett inn $x = 2$ i $f'(x)$.\\
Tolkning: & Grafen har stigningstall $4$ i punktet der $x = 2$.\\
\end{stegtabell}
\end{eksempel}

\begin{eksempel}[2 — gang ut først]
La $g(x) = (2x - 1)^2$. Finn $g'(1)$.

\begin{stegtabell}
$g(x) = 4x^2 - 4x + 1$ & Gang ut med kvadratsetningen. Vi har ingen regel for potens av parentes.\\
$g'(x) = 8x - 4$ & Deriver ledd for ledd.\\
$g'(1) = 8 - 4 = 4$ & Sett inn $x = 1$.\\
Kontroll: & $\dfrac{g(1+h) - g(1)}{h} = 4 + 4h$. For $h = 0,1$; $0,01$; $0,001$ får vi $4,4$; $4,04$; $4,004$. Tallene nærmer seg $4$.\\
\end{stegtabell}
\end{eksempel}

\begin{vanligfeil}
\begin{itemize}[leftmargin=1.2em,itemsep=1pt,topsep=0pt]
  \item $(x^3)' = 3x^3$ eller $3x^4$ er feil. Eksponenten ganges ned og blir én mindre: $(x^3)' = 3x^2$.
  \item $(5)' = 5$ og $(5x)' = 5x$ er feil. En konstant gir $0$, og $(5x)' = 5$.
  \item Den deriverte av et produkt er ikke produktet av de deriverte. Det gjelder ikke generelt: $(x^2\cdot x)' = (x^3)' = 3x^2$, men $(x^2)'\cdot(x)' = 2x\cdot 1 = 2x$.
  \item $s'(3)$ er ikke strekningen etter 3 sekunder. Strekningen er $s(3)$. Tallet $s'(3)$ er farten etter 3 sekunder.
\end{itemize}
\end{vanligfeil}

\nivaahode{1}

\oppg{1.1}{1}{\uh}
Deriver.
\begin{tasks}(4)
  \task $x^5$
  \task $4x^3$
  \task $6x - 9$
  \task $8$
\end{tasks}

\oppg{1.2}{1}{\uh}
Finn $f'(x)$.
\begin{tasks}(2)
  \task $f(x) = 3x^2 + 4x - 1$
  \task $f(x) = x^4 - 2x^3$
  \task $f(x) = \dfrac{1}{2}x^2 + 10$
  \task $f(x) = 5x^3 - x$
\end{tasks}

\oppg{1.3}{1}{\uh}
La $f(x) = x^3 - 3x^2 + 5x$. Fullfør.
\begin{tasks}(1)
  \task $f'(x) = \ \rute\, x^2 - \ \rute\, x + \ \rute$
  \task $f'(3) = \ \rute\cdot 3^2 - \ \rute\cdot 3 + \ \rute = $ \skrivlinje[2cm]
\end{tasks}

\nivaahode{2}

\oppg{1.4}{2}{\uh}
Gang ut eller forkort først. Deriver deretter. Skriv utregningen.
\begin{tasks}(2)
  \task $(x + 3)^2$
  \task $(3x - 1)(x + 2)$
  \task $\dfrac{x^4 + 2x^2}{x^2}$ \ ($x \neq 0$)
  \task $\dfrac{2x^3 - 6x^2}{2x}$ \ ($x \neq 0$)
\end{tasks}

\oppg{1.5}{2}{\uh}
Finn \textbf{den første feilen}. Løs riktig.
\[ f(x) = (x-1)(x+4) \ \longrightarrow\ f(x) = x^2 + 3x - 4 \ \longrightarrow\ f'(x) = 2x + 3 - 4 \ \longrightarrow\ f'(x) = 2x - 1 \ \longrightarrow\ f'(3) = 5 \]

\oppg{1.6}{2}{\mh}
En syklist har kommet $s(t) = 0,5t^2 + 2t$ meter etter $t$ sekunder (tenkt modell).
\begin{tasks}(1)
  \task Finn $s'(t)$.
  \task Regn ut $s'(4)$. Forklar hva tallet betyr, med enhet.
  \task Regn ut $s(4)$. Hva er forskjellen på $s(4)$ og $s'(4)$?
\end{tasks}

\nivaahode{3}

\oppg{1.7}{3}{\uh}
Er påstanden alltid, noen ganger eller aldri sann? Begrunn svaret.
\begin{tasks}(1)
  \task Hvis $f'(x) = g'(x)$ for alle $x$, er $f(x) = g(x)$ for alle $x$.
  \task Hvis $f$ har grad $n \ge 1$, har $f'$ grad $n - 1$.
\end{tasks}

\oppg{1.8}{3}{\uh}
Ola deriverer $(x+1)(x-1)$ faktor for faktor og får $1\cdot 1 = 1$.
\begin{tasks}(1)
  \task Deriver $(x+1)(x-1)$ riktig.
  \task For hvilken $x$ gir Olas metode likevel riktig stigningstall?
  \task Forklar hvorfor regelen til Ola ikke gjelder generelt.
\end{tasks}

\begin{sjekk}
\begin{tasks}(2)
  \task Deriver $3x^4 - 5x$.
  \task Hva er $(7)'$?
  \task! Hvorfor ganger vi ut $(x+2)^2$ før vi deriverer? Gjør det og deriver.
\end{tasks}
\end{sjekk}

% ---------------------------------------------------------------------

\del{2}{Tangentlikning}

\begin{ordliste}
tangentlikning & Likningen $y = mx + c$ til tangenten i et punkt på grafen. & \mittsprak\\
ettpunktsformelen & Likningen $y - y_1 = m(x - x_1)$ for en linje med stigningstall $m$ gjennom $(x_1, y_1)$. & \mittsprak\\
parallell & To linjer er parallelle hvis de har samme stigningstall. & \mittsprak\\
vannrett tangent & En tangent med stigningstall $0$. Den har likningen $y = $ et tall. & \mittsprak\\
\end{ordliste}

\begin{mustkunne}
\begin{minipage}[t]{0.56\linewidth}
\textbf{Tangenten i $x = a$} har
\begin{itemize}[leftmargin=1.2em,itemsep=1pt,topsep=2pt]
  \item stigningstall $m = f'(a)$, et \emph{tall}
  \item punktet $(a, f(a))$ felles med grafen
\end{itemize}
\textbf{Hvorfor:} En linje er bestemt av ett punkt og et stigningstall. Da gir ettpunktsformelen
\[ y - f(a) = f'(a)\,(x - a) \]
Regn om til $y = mx + c$.\par
\smallskip
\textbf{Gitt stigningstall $k$:} løs $f'(x) = k$. Vannrett tangent betyr $k = 0$. Parallell med $y = 3x + 1$ betyr $k = 3$.
\end{minipage}\hfill
\begin{minipage}[t]{0.40\linewidth}
\vspace{0pt}\centering
\begin{tikzpicture}[baseline=(current bounding box.north)]
\begin{axis}[width=\linewidth, height=0.95\linewidth, axis lines=middle, xmin=-0.6, xmax=4.0, ymin=-3.8, ymax=3.2,
  xtick={1}, xticklabels={$a$}, ytick={-2}, yticklabels={$f(a)$},
  tick label style={font=\scriptsize}, xlabel={$x$}, ylabel={$y$}, label style={font=\small}]
\addplot[domain=-0.4:3.5, thick, black!60, samples=40] {x^2 - 3*x};
\addplot[domain=-0.5:3.2, thick, hoved] {-x - 1};
\addplot[only marks, mark=*, mark size=1.5pt, hoved] coordinates {(1,-2)};
\node[font=\scriptsize, hoved, anchor=north east] at (axis cs:2.1,-3.1) {tangent};
\end{axis}
\end{tikzpicture}
\end{minipage}
\end{mustkunne}

\begin{eksempel}[1 — tangent i et punkt]
Finn tangenten til $f(x) = x^2 - 3x$ i $x = 1$.

\begin{stegtabell}
$f(1) = 1 - 3 = -2$ & Punktet på grafen er $(1, -2)$.\\
$f'(x) = 2x - 3$, så $f'(1) = -1$ & Stigningstallet er tallet $f'(1)$.\\
$y - (-2) = -1\cdot(x - 1)$ & Ettpunktsformelen.\\
$y + 2 = -x + 1$, så $y = -x - 1$ & Gang ut og løs for $y$.\\
Kontroll: & Punktet $(1, -2)$ gir $-1 - 1 = -2$. Og $x^2 - 3x = -x - 1$ gir $(x-1)^2 = 0$, så linja rører grafen bare i $x = 1$.\\
\end{stegtabell}
\end{eksempel}

\begin{eksempel}[2 — vannrett tangent]
Finn punktene der grafen til $f(x) = x^3 - 3x$ har vannrett tangent.

\begin{stegtabell}
$f'(x) = 3x^2 - 3$ & Deriver.\\
$3x^2 - 3 = 0$ & Vannrett tangent betyr stigningstall $0$.\\
$x^2 = 1$, så $x = -1$ eller $x = 1$ & Løs likningen.\\
$f(1) = -2$ og $f(-1) = 2$ & Regn ut $y$-verdiene. Punktene er $(1, -2)$ og $(-1, 2)$.\\
Tangentene er $y = -2$ og $y = 2$ & Vannrett linje gjennom punktet.\\
\end{stegtabell}
\end{eksempel}

\begin{vanligfeil}
\begin{itemize}[leftmargin=1.2em,itemsep=1pt,topsep=0pt]
  \item $y - f(a) = f'(x)(x - a)$ er feil. Stigningstallet må være tallet $f'(a)$. Med $f'(x) = 2x - 3$ får du en parabel, ikke en linje.
  \item Å bruke $f(a)$ som stigningstall er feil. For $f(x) = x^2 - 3x$ i $x = 1$ er $f(1) = -2$, men stigningstallet er $f'(1) = -1$.
  \item Å svare $x = a$ som vannrett tangent er feil. Likningen $f'(x) = 0$ gir bare $x$. Tangenten er $y = f(a)$.
  \item $y - (-4) = y - 4$ er feil. To minus gir pluss: $y - (-4) = y + 4$.
\end{itemize}
\end{vanligfeil}

\nivaahode{1}

\oppg{2.1}{1}{\uh}
Finn tangenten til $f(x) = x^2 + 3x$ i $x = 1$. Fullfør.
\begin{tasks}(1)
  \task $f(1) =$ \skrivlinje[2cm] \quad og \quad $f'(x) =$ \skrivlinje[3cm] \quad og \quad $f'(1) =$ \skrivlinje[2cm]
  \task $y - \ \rute\ = \ \rute\,(x - 1)$, så $y =$ \skrivlinje[3cm]
\end{tasks}

\oppg{2.2}{1}{\uh}
Finn tangenten på formen $y = mx + c$.
\begin{tasks}(1)
  \task $f(x) = x^2 - 4x$ i $x = 3$
  \task $g(x) = x^3$ i $x = 1$
  \task $h(x) = 2x^2 - x + 1$ i $x = 0$
\end{tasks}

\oppg{2.3}{1}{\uh}
La $f(x) = x^2 - 6x$. Finn $x$ slik at tangenten
\begin{tasks}(3)
  \task er vannrett
  \task har stigningstall $4$
  \task har stigningstall $-2$
\end{tasks}

\nivaahode{2}

\oppg{2.4}{2}{\uh}
La $f(x) = x^3 - 3x^2 + 2$.
\begin{tasks}(1)
  \task Finn tangenten i $x = 1$.
  \task Finn punktene på grafen der tangenten er vannrett.
  \task Finn alle punkter der tangenten er parallell med $y = 9x + 1$.
\end{tasks}

\oppg{2.5}{2}{\uh}
Finn \textbf{den første feilen}. Løs riktig.
\[ f(x) = x^2 - 5x \ \longrightarrow\ f(4) = -4 \ \longrightarrow\ f'(x) = 2x - 5 \ \longrightarrow\ f'(4) = 3 \ \longrightarrow\ y - 4 = 3(x - 4) \ \longrightarrow\ y = 3x - 8 \]

\oppg{2.6}{2}{\uh}
En ball kastes rett opp. Høyden er $h(t) = 20t - 5t^2$ meter etter $t$ sekunder (tenkt modell).
\begin{tasks}(1)
  \task Finn $h'(t)$. Hva er farten etter 1 sekund? Husk enhet.
  \task Når er farten $0$ m/s? Hva er høyden da?
  \task Finn toppunktet til $h$ med symmetrilinja fra 3.2. Sammenlign med b). Hvorfor må svarene stemme?
\end{tasks}

\nivaahode{3}

\oppg{2.7}{3}{\uh}
Parabelen $f(x) = ax^2 + bx$ har tangent med stigningstall $4$ i $x = 1$, og går gjennom punktet $(3, 15)$.
\begin{tasks}(1)
  \task Finn $a$ og $b$.
  \task Finn tangentlikningen i $x = 1$.
\end{tasks}

\oppg{2.8}{3}{\uh}
Er påstanden alltid, noen ganger eller aldri sann? Begrunn svaret.
\begin{tasks}(1)
  \task To ulike punkter på grafen til et andregradspolynom har parallelle tangenter.
  \task Grafen til et tredjegradspolynom har en vannrett tangent.
  \task Hvor mange vannrette tangenter kan grafen til et tredjegradspolynom ha? Begrunn.
\end{tasks}

\begin{sjekk}
\begin{tasks}(2)
  \task Finn tangenten til $f(x) = x^2$ i $x = 2$.
  \task Hvor er tangenten til $f(x) = x^2 - 8x$ vannrett?
  \task! Hvorfor bruker vi tallet $f'(a)$ og ikke funksjonen $f'(x)$ som stigningstall i tangentlikningen?
\end{tasks}
\end{sjekk}
